// MALLOC
Ele cria um espaço na memória  
Ex:
    Elem* no = (Elem*) malloc(sizeof(Elem));
    Crie um novo nó na memória e faça o ponteiro no apontar para ele.
depois do malloc sempre se coloca essa verificação:
    if(no == NULL)
        return 0;



// FREE
AQUI SEMPRE COLOQUE SEM OS SINAIS DE PONTEIRO, POIS O FREE RECEBE O ENDEREÇO DO ESPAÇO A SER LIBERADO, E NÃO O VALOR DO PONTEIRO.
Libera o espaço alocado na memória
Ex: 
    free(no); // libera um no que nao sera mais usado
Exemplo de utilização:
    Elem* no = pi->topo; // salva o nó que será removido
    pi->topo = pi->topo->prox; // move o topo para o próximo
    free(no); // libera o nó antigo

    "salva em aux/no → move o ponteiro principal → free(aux/no)"



// .C e .H
Tudo que eh tipo de dado(struct, typedef struct elemento Pilha, enum) vai pro .h e implemetacoes e elementos(struct elemento, typedef strct elemento Elem) vao pro .c



// PILHA:
CRIA PILHA
Pilha* cria_Pilha(){
    Pilha* pi = (Pilha*) malloc(sizeof(Pilha));
    if(pi != NULL)
        *pi = NULL;
    return pi;
}

Pilha* pi = ponteiro para a estrutura da pilha.
pi->topo = ponteiro para o primeiro nó da pilha. esse so funciona se tiver aa struct pilha
ou 
*pi que tambem aponta pro topo
Elem* no = ponteiro para um nó da pilha.
no->prox = ponteiro para o próximo nó.

FALA SE A PILHA EXISTE
if(pi == NULL)
    return 0;

FALA SE TEM ALGUM ELEMENTO NA PILHA
if(pi->topo == NULL) // Ela existe mas nao tem nada dentro
    return 0;


CRIA NO:
Elem* no; // cria ponteiro para o novo nó

no = (Elem*) malloc(sizeof(Elem)); // reserva memória para o nó

if(no == NULL) // verifica se o malloc falhou
    return 0;


REMOVER SEM PERDER PONTEIRO:
Elem* no = pi->topo; // salva o nó que será removido

pi->topo = pi->topo->prox; // topo anda para o próximo

free(no); // libera o antigo topo



// FILA:

